\documentclass[11pt]{article}
\usepackage{amsmath,amssymb,amsthm,algorithm,algorithmic}
\usepackage[margin=1in]{geometry}

\newtheorem{theorem}{Theorem}
\newtheorem{lemma}[theorem]{Lemma}
\newtheorem{corollary}[theorem]{Corollary}

\title{Problem 6: Sum Square Difference}
\author{Project Euler Solutions}
\date{}

\begin{document}
\maketitle

\section{Problem Statement}

Compute the difference between the square of the sum and the sum of the squares of the first 100 natural numbers:
\[
D(100) = \left(\sum_{i=1}^{100} i\right)^{\!2} - \sum_{i=1}^{100} i^2.
\]

\section{Notation}

Throughout, write
\[
S_1(n) = \sum_{i=1}^{n} i, \qquad S_2(n) = \sum_{i=1}^{n} i^2, \qquad D(n) = S_1(n)^2 - S_2(n).
\]

\section{Mathematical Development}

\begin{lemma}[Sum of Natural Numbers]\label{lem:s1}
For all $n \ge 1$, \; $S_1(n) = n(n+1)/2$.
\end{lemma}

\begin{proof}
By induction. Base: $S_1(1) = 1 = 1 \cdot 2/2$. Step: assuming $S_1(k) = k(k+1)/2$,
\[
S_1(k+1) = \frac{k(k+1)}{2} + (k+1) = \frac{(k+1)(k+2)}{2}. \qedhere
\]
\end{proof}

\begin{lemma}[Sum of Squares]\label{lem:s2}
For all $n \ge 1$, \; $S_2(n) = n(n+1)(2n+1)/6$.
\end{lemma}

\begin{proof}
The telescoping identity $(k+1)^3 - k^3 = 3k^2 + 3k + 1$, summed from $k=1$ to $n$, gives
\[
(n+1)^3 - 1 = 3S_2(n) + 3S_1(n) + n.
\]
Substituting $S_1(n) = n(n+1)/2$ from Lemma~\ref{lem:s1} and solving:
\begin{align*}
3S_2(n) &= (n+1)^3 - 1 - \tfrac{3n(n+1)}{2} - n \\
        &= \frac{2n^3 + 3n^2 + n}{2} = \frac{n(n+1)(2n+1)}{2},
\end{align*}
where the last equality uses $2n^2 + 3n + 1 = (2n+1)(n+1)$. Dividing by 3 yields the result.
\end{proof}

\begin{theorem}[Closed Form]\label{thm:closed}
For all $n \ge 1$,
\[
D(n) = \frac{n(n-1)(n+1)(3n+2)}{12}.
\]
\end{theorem}

\begin{proof}
By Lemmas~\ref{lem:s1} and~\ref{lem:s2},
\[
D(n) = \frac{n^2(n+1)^2}{4} - \frac{n(n+1)(2n+1)}{6} = \frac{n(n+1)}{12}\bigl[3n(n+1) - 2(2n+1)\bigr].
\]
Expanding the bracket:
\[
3n(n+1) - 2(2n+1) = 3n^2 + 3n - 4n - 2 = 3n^2 - n - 2.
\]
The discriminant of $3n^2 - n - 2$ is $\Delta = 1 + 24 = 25$, giving roots $n = 1$ and $n = -2/3$. Hence
\[
3n^2 - n - 2 = (n - 1)(3n + 2),
\]
and $D(n) = n(n+1)(n-1)(3n+2)/12$.
\end{proof}

\begin{corollary}[Non-negativity]
$D(n) \ge 0$ for all $n \ge 1$, with equality iff $n = 1$.
\end{corollary}

\begin{proof}
For $n \ge 2$, each factor $n, n-1, n+1, 3n+2$ is positive. For $n = 1$: $D(1) = 0$.
\end{proof}

\begin{theorem}[Combinatorial Interpretation]\label{thm:comb}
$\displaystyle D(n) = 2\sum_{1 \le i < j \le n} ij$.
\end{theorem}

\begin{proof}
Expanding $\bigl(\sum_{i=1}^n i\bigr)^2 = \sum_i i^2 + 2\sum_{i<j} ij$ (by separating diagonal and off-diagonal terms in the double sum $\sum_{i,j} ij$) and subtracting $\sum i^2$ yields the result.
\end{proof}

\section{Algorithm}

\begin{algorithm}[H]
\caption{Sum square difference via closed form}
\begin{algorithmic}[1]
\RETURN $n \cdot (n-1) \cdot (n+1) \cdot (3n+2) \;/\; 12$
\end{algorithmic}
\end{algorithm}

\textbf{Application to $n = 100$:}
\[
D(100) = \frac{100 \times 99 \times 101 \times 302}{12} = \frac{301\,969\,800}{12} = 25\,164\,150.
\]

\emph{Verification.} $S_1(100) = 5050$, $S_1(100)^2 = 25\,502\,500$, $S_2(100) = 338\,350$, and $25\,502\,500 - 338\,350 = 25\,164\,150$.\;\checkmark

\section{Complexity Analysis}

\begin{theorem}[Complexity]
The closed-form algorithm runs in $O(1)$ time and $O(1)$ space.
\end{theorem}

\begin{proof}
The formula requires exactly 4 multiplications and 1 division. A brute-force approach iterating to compute $S_1$ and $S_2$ would require $O(n)$ time.
\end{proof}

\section{Answer}

\[
\boxed{25164150}
\]

\end{document}
